% !PRINTED 477
\weberpage{477}

%<<CONT
and by \S.\,146, 1. we thus have, for every index $k$,
\weq{2}{\varrho_k = \Phi[\chi(\varrho_k), \chi_1(\varrho_k) \ldots \chi_{m-1}(\varrho_k)].}
Now if the substitutions $\sigma_h, \sigma_k$ produce the same permutation among the $\alpha$, then for every index $s = 0, 1 \ldots m-1$
\[
\chi_s(\varrho_h) = \chi_s(\varrho_k),
\]
and from this it follows by (2) that we must also have $\sigma_h = \sigma_k$.

If the substitution $\sigma_h$ produces the permutation $\pi_h$ among the roots of $F(x) = 0$, then we say that \gesperrt{the permutation $\pi_h$ corresponds to the substitution $\sigma_h$}.

To every substitution there corresponds a definite permutation, and the number of these permutations — counting in the identical permutation, which corresponds to the identical substitution — is therefore equal to $\mu$.

If the permutations $\pi_h$ and $\pi_k$ correspond to the substitutions $\sigma_h$ and $\sigma_k$, then to the composed substitution $\sigma_h \sigma_k$ there corresponds the composed permutation $\pi_h \pi_k$. For the permutations are the result of the corresponding substitutions on certain quantities of the field $\Omega(\varrho)$, and here it makes no difference whether we carry out two substitutions one after the other or carry out the substitution composed of both at one stroke (by \S.\,147). We therefore have the theorem:

\gesperrt{I. To the $\mu$ substitutions of the field $\Omega(\varrho)$}
\weq{3}{\sigma, \sigma_1, \sigma_2 \ldots \sigma_{\mu-1},}
\gesperrt{there correspond $\mu$ permutations of $m$ figures}
\weq{4}{\pi, \pi_1, \pi_2 \ldots \pi_{\mu-1}}
\gesperrt{in such a way that to the substitution composed of two substitutions there corresponds the permutation composed of the corresponding permutations.}

But from this it follows that the permutations (4) form a group. We denote this group by $P$ and call it the \gesperrt{Galois group of the equation $F(x) = 0$, or also the Galois group of each of the fields $\Omega(\alpha), \Omega(\alpha_1) \ldots \Omega(\alpha_{m-1})$}.

This group stands to the group of substitutions of the field $\Omega(\varrho)$ in the relation expressed by Theorem I., for which reason the two groups are called \gesperrt{isomorphic}.
